%! Function Fundamentals: Notation, Domain, and Range — Unit 2, Lesson 2.1 — Group Activity (Answer Key)
\documentclass[10pt]{article}
\usepackage{atda-article}
\usepackage{atda-key}   % pulls in -boxes; do NOT also load -boxes
\newcommand{\TallMath}[1]{$\displaystyle #1\rule[-1.4em]{0pt}{3.2em}$}
\setlength{\workrowsep}{6pt}

\begin{document}

\pageheader{Unit 2, Lesson 2.1}{Group Activity --- Answer Key}
% NAMESTRIP: no \namedateperiod here — mirrors the blank. Teacher-only prose goes
% in the lesson plan as \begin{teachernote}[Group Activity], never in a key.

\vspace{0.15in}

\begin{scenariobox}[The Scenario]{royal}
\small
Every graph below gets asked \emph{both} questions: given an input, what is the output --- and
given an output, which input produced it? Every answer also gets written as a point. Work with
your group. \textbf{Everyone starts at Tier R.} When your whole group can do Tier R without help,
move on.
\end{scenariobox}

\vspace{0.10in}

\boxguard[30]
\begin{tcolorbox}[colback=white, colframe=black!40, title=\textbf{Tier R --- Remediate},
    fonttitle=\bfseries, arc=2mm, left=3mm, right=3mm, top=2mm, bottom=2mm, breakable]
\small
A candle is lit and burns at a steady rate. $L$ is its length in centimeters, $t$ hours after it
was lit.

\begin{center}
\begin{tikzpicture}
\begin{axis}[
    width=9.6cm, height=4.6cm,
    xlabel={$t$ (hours since lighting)}, ylabel={$L$ (length, cm)},
    xmin=0, xmax=10.6, ymin=0, ymax=22,
    xtick={0,2,4,6,8,10}, minor xtick={1,3,5,7,9}, ytick={0,4,8,12,16,20},
    grid=both, grid style={linegray!45}, axis lines=left,
    tick label style={font=\scriptsize}, label style={font=\scriptsize}]
  \addplot[royal, very thick, domain=0:10, samples=2]{20-2*x};
\end{axis}
\end{tikzpicture}
\end{center}

\begin{enumerate}[leftmargin=1.6em, itemsep=4pt, topsep=2pt]
  \item Read $L(6)$ off the graph. Write it as a point, then as a sentence about the candle.

\ansline{$L(6)=8$, so $(6,8)$ is on the graph.}\\
\ansline{Six hours after it was lit the candle is 8 cm long.}\\

  \item Now the other direction: for what value of $t$ is $L(t)=4$? Read it off the graph, then
        check it against the formula $L(t)=20-2t$.
\begin{work}
  20-2t &= 4 \\
   -2t &= -16 \\
     t &= 8
\end{work}

  \item State the domain and the range in this situation, with units, in inequality notation
        \emph{and} in interval notation.

\ansline{Domain: $0 \le t \le 10$ hours, or $[0,10]$.}\\
\ansline{Range: $0 \le L \le 20$ centimeters, or $[0,20]$.}\\

  \item Why does the domain stop at $t=10$? The formula would happily give $L(12)=-4$.

\ansline{At $t=10$ the candle is gone; a candle cannot have a negative length.}\\
\end{enumerate}
\end{tcolorbox}

\vspace{0.10in}

\boxguard[30]
\begin{tcolorbox}[colback=white, colframe=black!40,
    title=\textbf{Tier A --- Approaching Proficiency},
    fonttitle=\bfseries, arc=2mm, left=3mm, right=3mm, top=2mm, bottom=2mm, breakable]
\small
A skate park bowl is $12$ feet across. $H$ is the height of the concrete surface in feet above the
lowest point of the bowl, measured $x$ feet in from the left rim.

\begin{center}
\begin{tikzpicture}
\begin{axis}[
    width=10.4cm, height=4.8cm,
    xlabel={$x$ (feet in from the left rim)}, ylabel={$H$ (height, feet)},
    xmin=0, xmax=12.6, ymin=0, ymax=13.5,
    xtick={0,2,4,6,8,10,12}, minor xtick={1,3,5,7,9,11}, ytick={0,3,6,9,12},
    grid=both, grid style={linegray!45}, axis lines=left,
    tick label style={font=\scriptsize}, label style={font=\scriptsize}]
  \addplot[royal, very thick, domain=0:12, samples=80]{(x-6)^2/3};
\end{axis}
\end{tikzpicture}
\end{center}

\begin{enumerate}[leftmargin=1.6em, itemsep=4pt, topsep=2pt]
  \item Read $H(3)$ off the graph. Write it as a point, then as a sentence about the bowl.

\ansline{$H(3)=3$, so $(3,3)$ is on the graph.}\\
\ansline{Three feet in from the left rim, the concrete is 3 feet above the bottom.}\\

  \item Find \emph{every} value of $x$ for which $H(x)=3$. Say what those two answers mean for
        somebody standing in the bowl.

\ansline{$x=3$ and $x=9$.}\\
\ansline{The bowl is symmetric: the same height occurs on the left wall and on the right wall.}\\

  \item State the domain and the range, with units, in inequality notation \emph{and} in interval
        notation.

\ansline{Domain: $0 \le x \le 12$ feet, or $[0,12]$.}\\
\ansline{Range: $0 \le H \le 12$ feet, or $[0,12]$.}\\

  \item Exactly one height in the range comes from exactly one value of $x$. Which height, which
        $x$, and why is it the only one?

\ansline{$H=0$ at $x=6$ only --- the single lowest point of the bowl.}\\
\ansline{Every other height is reached once on each wall, so it has two inputs.}\\

  \item A contractor asks what $H(15)$ is. What do you tell them, and why is that not a question
        about arithmetic?

\ansline{There is no $H(15)$: $15$ is outside the domain, because the bowl is only 12 ft across.}\\
\end{enumerate}
\end{tcolorbox}

\vspace{0.10in}

\boxguard[30]
\begin{tcolorbox}[colback=white, colframe=black!40, title=\textbf{Tier E --- Extension},
    fonttitle=\bfseries, arc=2mm, left=3mm, right=3mm, top=2mm, bottom=2mm, breakable]
\small

\begin{enumerate}[leftmargin=1.6em, itemsep=4pt, topsep=2pt]
  \item \textbf{Critique the work.} Two students wrote these about the Tier R candle. Each is
        wrong. Name the error and give the correct answer.
        \begin{enumerate}[label=(\alph*), leftmargin=1.6em, itemsep=3pt, topsep=3pt]
          \item ``$L(6)=8$, so $L(8)=6$.''

\ansline{Function notation does not flip. The input goes inside the parentheses,}\\
\ansline{the output comes after the equals sign. In fact $L(8)=20-2(8)=4$.}\\

          \item ``The range is $[20,0]$, because the candle starts at $20$ cm and ends at $0$ cm.''

\ansline{An interval is always written smaller number first: the range is $[0,20]$.}\\
        \end{enumerate}

  \item \textbf{A graph made of dots.} At the concession stand a hot dog costs \$3, and a family
        may buy at most $8$. $P(n)$ is the price paid for $n$ hot dogs.

\begin{center}
\begin{tikzpicture}
\begin{axis}[
    width=8.6cm, height=4.0cm,
    xlabel={$n$ (hot dogs)}, ylabel={$P$ (price, dollars)},
    xmin=-0.4, xmax=8.6, ymin=0, ymax=27,
    xtick={0,1,2,3,4,5,6,7,8}, ytick={0,6,12,18,24},
    grid=both, grid style={linegray!45}, axis lines=left,
    tick label style={font=\scriptsize}, label style={font=\scriptsize}]
  \addplot[royal, only marks, mark=*, mark size=1.6pt]
    coordinates {(0,0) (1,3) (2,6) (3,9) (4,12) (5,15) (6,18) (7,21) (8,24)};
\end{axis}
\end{tikzpicture}
\end{center}

        \begin{enumerate}[label=(\alph*), leftmargin=1.6em, itemsep=3pt, topsep=3pt]
          \item Write the domain and the range of $P$ in set notation.

\ansline{Domain: $\{0, 1, 2, 3, 4, 5, 6, 7, 8\}$ hot dogs.}\\
\ansline{Range: $\{0, 3, 6, 9, 12, 15, 18, 21, 24\}$ dollars.}\\

          \item A student writes ``domain $[0,8]$, range $[0,24]$.'' Give one number that their
                answer allows and the situation does not, and say why interval notation is the
                wrong tool here.

\ansline{$[0,8]$ allows $2.5$ hot dogs, and $[0,24]$ allows a price of \$7.}\\
\ansline{An interval holds every number between its ends; this domain is a list of dots.}\\
        \end{enumerate}

  \item \textbf{Justify.} A classmate says: ``The skate bowl graph breaks the definition of a
        function, because the height $3$ feet happens at two different places.'' Are they right?
        Use the bowl to explain what the definition actually promises.

\ansline{No. A function promises each \emph{input} exactly one output --- not the reverse.}\\
\ansline{At $x=3$ the concrete is at exactly one height, and the same at $x=9$.}\\
\ansline{Two inputs are allowed to share an output; one input may never have two.}\\
\end{enumerate}
\end{tcolorbox}

\end{document}
